% ===================================================================
%  TABEL Nr. 1 — Kool. --- Gümnaasium. --- Klassiruum.
%
%  Ceci n'est PAS une transcription : c'est une traduction, faite en
%  2026, en estonien standard d'aujourd'hui. D'ou trois differences
%  avec texto/io et texto/fr, et trois seulement :
%
%   * pas de \nl ni de \cc — la traduction ne suit aucune ligne
%     imprimee, et les lignes de ce fichier ne sont que du confort ;
%   * pas de \begin{VUpage} — elle n'a ni feuillet ni folio, et la
%     page de lecture ne lui en affiche pas ;
%   * le gras se pose sur le TERME seul, ponctuation dehors.
%
%  Tout le reste est identique, et doit l'etre : les cles « %%K » sont
%  celles de texto/io, macro pour macro pour l'apparat, et les renvois
%  \textsuperscript{(n)} visent les MEMES objets de la planche.
%
%  LA DIX-SEPTIEME ET DERNIERE DES LANGUES DE LA COMMUNAUTE IDISTE.
%  L'ido a eu ses groupes a Tallinn et a Tartu entre les deux guerres.
%
%  UNE DECISION DE COMPOSITION, ECRITE ICI POUR QU'ON NE LA CHERCHE
%  PAS AILLEURS. L'estonien ecrit « tabel », qui est un morceau du
%  « TABELO » de l'ido : le motif NUMERO_TAB de outils/html.py ne peut
%  donc pas le prendre seul, et exige le « Nr » qu'il exige deja du
%  neerlandais et du suedois. La colonne compose donc « TABEL Nr. 1 »
%  avec un N capital, comme ces deux-la. Le cablage n'a rien recu de
%  ce cote ; il a recu les quatre ordinaux, « seeria » et « stseen ».
%
%  QUELLE LANGUE ECRIRE ? LA QUESTION SE POSE ICI PLUS QU'AILLEURS.
%  Le livret est de 1926, c'est-a-dire en plein milieu de la
%  « keeleuuendus » de Johannes Aavik : la seule entreprise en Europe
%  ou l'on ait FABRIQUE des racines de toutes pieces — « relv »
%  l'arme, « roim » le crime, « laip » le cadavre, « veenma »
%  convaincre, « ese » l'objet — sans les tirer d'aucune langue.
%  L'usage en a garde une partie et rejete le reste.
%
%  LA COLONNE EST ECRITE EN ESTONIEN STANDARD D'AUJOURD'HUI, celui
%  qui a fait le tri. Le romanche du tableau 1 precedent avait une
%  ORTHOGRAPHE inventee en 1982 pour des mots qui existaient ;
%  l'estonien a des MOTS inventes vers 1920 pour une orthographe qui
%  existait. Les deux dernieres colonnes de la serie sont donc les
%  deux seules ou une decision d'homme est passee dans la langue —
%  et ce n'est pas la meme moitie de la langue.
%
%  ET LA PAGE DONNE DEJA UNE CARTE. Le mobilier de la classe vient de
%  l'ouest, du bas-allemand des villes hanseatiques : « tahvel » le
%  tableau, « kriit » la craie, « pliiats » le crayon, « sirkel » le
%  compas, « paber » le papier, « kool » l'ecole meme. Mais LE LIVRE
%  VIENT DE L'EST : « raamat » est le « gramota » russe, le mot grec
%  des lettres passe par Novgorod. Les objets sont arrives par les
%  marchands, l'ecrit par les moines, et l'estonien porte les deux
%  routes dans la meme salle.
%
%  LES NOMS DE PERSONNE SONT ESTONIENS : Ioannes est Jaan, Petrus est
%  Peeter, Karolus est Kaarel, Georgius est Jüri, Iakobus est Jaakob.
%  Les noms latins de la note (*) ne bougent pas : c'est d'eux que la
%  note parle.
% ===================================================================

%%K t01-apar-0 apar
\VUsaut{2.0mm}
\VUcentre{12.6pt}{120}{SELETUSVIHIK}
\VUsaut{3.2mm}
\VUcentre{8.4pt}{160}{TEOSELE}
\VUsaut{2.6mm}
\VUtitre{15.6pt}{0}{Delmas'i abitabelid}
\VUsaut{3.4mm}
\VUornamento{9pt}
\VUsaut{5.0mm}
\VUcentre{13.2pt}{90}{ESIMENE SEERIA}
\VUsaut{3.0mm}
\VUfilet{16mm}
\VUsaut{5.6mm}

%%K t01-apar-1 apar
\VUtitre{15.0pt}{60}{TABEL Nr. 1}
\VUsaut{1.4mm}
\VUtitre{11.4pt}{0}{Kool. --- Gümnaasium. --- Klassiruum.}
\VUsaut{2.2mm}
\VUfilet{20mm}
\VUsaut{6.4mm}

%%K t01-tit-1 sub
\VUpk{10.2pt}{40}{Üldine kirjeldus.}
\VUsaut{3.0mm}

%%K t01-01-1 p
1. --- See tabel kujutab \VUgras{gümnaasiumi} \VUgras{klassiruumi}. Selles
klassiruumis on kaheksa \VUgras{lauda} \textsuperscript{(18)} ja kaheksa
\VUgras{pinki} \textsuperscript{(32)}. Igal pingil istub kolm õpilast.
\VUgras{Õpetaja} \textsuperscript{(7)} istub \VUgras{toolil} \textsuperscript{(8)} oma
\VUgras{kateedri} \textsuperscript{(6)} taga.

%%K t01-02-1 p
2. --- Kateedrist vasakul on \VUgras{kapp} klaasustega \textsuperscript{(15)}. Paremal
on \VUgras{must tahvel} \textsuperscript{(1)} koos \VUgras{lapiga} \textsuperscript{(2)} ja
\VUgras{kriidiga} \textsuperscript{(3)}.

%%K t01-02-2 p
Kateedri kohal ripuvad \VUgras{seinatahvlid} \textsuperscript{(9, 11, 12)} ja
\VUgras{kaart} \textsuperscript{(10)}.

%%K t01-03-1 p
3. --- Mitte kõik õpilased ei istu oma \VUgras{kohal} \textsuperscript{(17)}. Jaan
\textsuperscript{(4)} seisab tahvli juures; ta hoiab lappi ja kustutab
\VUgras{geomeetrilisi kujundeid}.

%%K t01-03-2 p
Peeter on kateedri lähedal; ta kogub \VUgras{kirjalikke töid} \textsuperscript{(5)} ja
viib need õpetajale.

%%K t01-04-1 p
4. --- \VUgras{See} õpetaja on \VUgras{elavate keelte} õpetaja. Ta õpetab õpilasi
rääkima, lugema ja kirjutama võõrkeeli \VUgras{(saksa, inglise, hispaania, itaalia,
vene} või \VUgras{prantsuse)}.

%%K t01-05-1 p
5. --- Klassiruum on väga avar ja väga valge. Tagaseinas on lai \VUgras{aken}, kaks
\VUgras{tuulutusakent} \textsuperscript{(78)} ja \VUgras{klaasuks} selle õhutamiseks ja
valgustamiseks.

%%K t01-05-2 p
See klassiruum ei ole külm. Keskel on selle soojendamiseks väga hea \VUgras{ahi}
\textsuperscript{(46)} oma \VUgras{toruga} \textsuperscript{(45)}.

%%K t01-05-3 p
Seina külge on kinnitatud \VUgras{riidenagid} \textsuperscript{(58)}; nende otsas
ripuvad õpilaste mütsid ja mantlid.

%%K t01-tit-2 sub
\VUsaut{5.2mm}
\VUpk{10.2pt}{40}{Mänguõu.}
\VUsaut{3.6mm}

%%K t01-06-1 p
6. --- \VUgras{Kell} \textsuperscript{(63)} näitab üheksa ja viis minutit.

%%K t01-06-2 p
\VUgras{Vahetund} \textsuperscript{(76)} on just lõppenud ja tund algab uuesti.
Vahetund peetakse \VUgras{õues} \textsuperscript{(75)}. Me näeme mõnda mängivat
õpilast. \VUgras{Kasvataja} \textsuperscript{(74)} valvab neid.

%%K t01-07-1 p
7. --- Õues näeme veel teisi inimesi: \VUgras{inspektorit} \textsuperscript{(73)},
\VUgras{direktorit} \textsuperscript{(71)} ja \VUgras{õppealajuhatajat}
\textsuperscript{(72)}.

%%K t01-07-2 p
Inspektor kontrollib koole, direktor juhib gümnaasiumi, õppealajuhataja valvab
õppetööd.

%%K t01-07-3 p
Neljas isik on \VUgras{uksehoidja} \textsuperscript{(77)}. Ta kannab kaenla all
päevikut, kuhu märkida puudujad.

%%K t01-08-1 p
8. --- Õue tagaosas on \VUgras{võimlemisriistad} \textsuperscript{(70)} ja
\VUgras{käsitöö} \textsuperscript{(69)} töökoda.

%%K t01-08-2 p
Tabelist paremal hakkame nägema \VUgras{muusika}- ja \VUgras{joonistusruume}.

%%K t01-noto-1 noto
\VUnotes{143.2mm}{%
(*) \textit{Ristinimede} puhul soovitatakse need kirjutada ka ladinapärasel kujul.
See on ainus viis vältida loendamatuid erikujusid. Kuidas muidu valida
\textit{Giovanni, Jean,} \textit{Johann, Jan, Hans, John, Ivan} jne. vahel? Kas me
loome ristinimede jaoks eraldi sõnaraamatu? Võtkem ladina keelest nende nimetav
kääne ja öelgem: \VUgras{Ioannes, Ioanna; Iulius, Iulia; Iakobus; Andreas; Lukas.}
(Täielik grammatika).%
}

%%K t01-tit-3 sub
\VUpk{10.2pt}{40}{Üksikasjalik kirjeldus.}
\VUsaut{2.4mm}

%%K t01-tit-4 sub
\VUpk{10.2pt}{40}{Head õpilased.}
\VUsaut{3.0mm}

%%K t01-09-1 p
9. --- Kateedrist paremal esimeses pingis istuvad \VUgras{Henrik} (*),
\VUgras{Stefan} ja \VUgras{Kaarel}.

%%K t01-09-2 p
Henrik on väga tähelepanelik ja töökas; ta kuulab õpetajat. Tema laual on
\VUgras{vihik} \textsuperscript{(28)}, \VUgras{raamat} \textsuperscript{(29)},
\VUgras{sõnaraamat} \textsuperscript{(30)} ja \VUgras{pinal} \textsuperscript{(31)}.
Tema raamatul on palju \VUgras{lehekülgi}; igas peatükis on mitu \VUgras{lõiku} ja
igal \VUgras{real} palju \VUgras{sõnu}.

%%K t01-09-4 p
Tema naaber Stefan \textsuperscript{(24)} on hoolas. Ta märgib oma vihikusse järgmise
korra õppetüki. Tal on ilus \VUgras{käekiri}. Tema raamat on suletud. Me näeme ainult
\VUgras{kaant} \textsuperscript{(27)}. Tema kõrval on \VUgras{sirkel}
\textsuperscript{(26)}.

%%K t01-09-5 p
Kaarel \textsuperscript{(23)} hoiab raamatut käes; ta avab selle, et õppetükki üle
korrata. Nagu igal õpilasel on temal enda ees \VUgras{tindipott}
\textsuperscript{(25)}. Ta on sõnakuulelik.

%%K t01-10-1 p
10. --- Kõrvallauas istuvad \VUgras{Ludvig, Anton} ja \VUgras{Paul}. Ludvig vastab
oma \VUgras{õppetükki} \textsuperscript{(22)} ja vastab õpetaja \VUgras{küsimustele}.
Anton \textsuperscript{(21)} on väga tähelepanematu; mõnikord õpetaja isegi karistab
teda. Paul \textsuperscript{(20)} on hea \VUgras{kaaslane}, lahke ja abivalmis. Ta on
väga tähelepanelik.

%%K t01-tit-5 sub
\VUsaut{4.2mm}
\VUpk{10.2pt}{40}{Halvad õpilased.}
\VUsaut{3.2mm}

%%K t01-11-1 p
11. --- Henriku taga istub \VUgras{Jüri} \textsuperscript{(38)}. Ta ei ole halb poiss,
kuid ta on \VUgras{lobisev}, tähelepanematu ja rahutu. Tema \VUgras{kergemeelsus} ja
\VUgras{sõnakuulmatus} tekitavad tunnis \VUgras{korratust} ja sageli teenib ta karmi
\VUgras{hoiatuse}. Tema \VUgras{mustandivihik} \textsuperscript{(35)} on räpane, ta
teeb sinna \VUgras{tindiplekke} oma \VUgras{sulepeaga} \textsuperscript{(34)}, ja
seepärast kasutab ta alati oma \VUgras{kustutuskummi} \textsuperscript{(36)}; tema
\VUgras{portfell} \textsuperscript{(33)} on rebenenud, sest ta on väga hooletu. Miks ta
toob oma \VUgras{pliiatsihoidja} \textsuperscript{(37)} elavate keelte tundi?

%%K t01-11-2 p
Tema naaber Edmund \textsuperscript{(39)} ei salli teda. Ta on tõesti pisut laisk,
kuid ta on vaikne ja püsib rahulikult. Ta ei lobise; ainult et kuulamise asemel
eelistab ta lugeda oma \VUgras{lugemiku} \VUgras{jutte}.

%%K t01-11-3 p
Selle pingi kolmas õpilane on \VUgras{Ludvig} \textsuperscript{(40)}. Ta on hoolas. Ta
kirjutab valgele \VUgras{paberilehele}. Enda ette on ta pannud oma
\VUgras{kuivatuspaberi} \textsuperscript{(41)}.

%%K t01-12-1 p
12. --- Teise pingi õpilased, õpetajast vasakul, on väga noored. Esimene, väike
\VUgras{Emil}, ei oska veel hästi kirjutada; ta toob kaasa oma
\VUgras{krihvlitahvli} \textsuperscript{(42)}, oma \VUgras{krihvli} ja oma
\VUgras{käsna} \textsuperscript{(43)}.

%%K t01-12-2 p
Tema kõrval istuvad \VUgras{August} ja \VUgras{Bernard} \textsuperscript{(44)}. See
viimane töötab hästi, kuid tema \VUgras{riietus} on väga hooletu.

%%K t01-13-1 p
13. --- Nende taga istub kolm \VUgras{laiska. Albert} ei õpi oma õppetükke.
\VUgras{Jaakob} \textsuperscript{(47)} magab kuulamise asemel. Tema \VUgras{laiskus}
toob talle \VUgras{etteheiteid} ja isegi \VUgras{karistusi}. Lõpuks on
\VUgras{Kristjan} \textsuperscript{(48)} liiga kergemeelne. Ta \VUgras{käitub} halvasti
ega tee mingit pingutust, et end parandada.

%%K t01-14-1 p
14. --- Tema naabrid seevastu käituvad hästi. \VUgras{Ernst} \textsuperscript{(51)}
armastab \VUgras{korda} ja \VUgras{korralikkust}; ta kasutab alati oma
\VUgras{joonlauda} \textsuperscript{(50)} ja oma \VUgras{pliiatsit}
\textsuperscript{(49)}. Ta on \VUgras{päevaõpilane}.

%%K t01-14-2 p
\VUgras{Frants} \textsuperscript{(52)} on \VUgras{internaadiõpilane}; ta kannab
gümnaasiumi \VUgras{vormiriietust}; ta sööb ja magab seal. Ta on hea
\VUgras{kaaslane}.

%%K t01-14-3 p
Moritz teritab oma \VUgras{pliiatsit} oma \VUgras{taskunoaga} \textsuperscript{(56)};
ta on väga hoolikas.

%%K t01-14-4 p
Ta toob kaasa kõik oma raamatud, oma \VUgras{grammatika} \textsuperscript{(55)}, oma
\VUgras{märkmiku} \textsuperscript{(54)} ja isegi \VUgras{kirjutusaluse}
\textsuperscript{(53)}. Tema \VUgras{koolikott} \textsuperscript{(57)} on põrandal tema
taga.

%%K t01-14-5 p
Klassiruumi tagumist kahte lauda hõivavad \VUgras{head õpilased}
\textsuperscript{(19)}.

%%K t01-tit-6 sub
\VUsaut{4.6mm}
\VUpk{10.2pt}{40}{Joonistamine ja muusika.}
\VUsaut{3.4mm}

%%K t01-15-1 p
15. --- \VUgras{Joonistusklassis} on ilusad \VUgras{mudelid} seina küljes ja
\VUgras{lamp} \textsuperscript{(62)} \VUgras{lambivarjuga}, mis ripub lae küljes.
\VUgras{Õpetaja} \textsuperscript{(60)} on väga kannatlik; ta parandab õpilaste
\VUgras{joonistusi} \textsuperscript{(61)} ja õpetab neid natuurist \VUgras{büsti}
\textsuperscript{(59)} joonistama.

%%K t01-16-1 p
16. --- \VUgras{Muusikaklass} tundub väga huvitav. Seal õpitakse \VUgras{muusikat}
lugema, solfedžeerima \VUgras{gamminoote} \textsuperscript{(67)}, mis on kirjutatud
\VUgras{noodijoonestikule} (do, re, mi, fa, sol, la, si\,=\,C. D. E. F. G. A. B.),
tundma \VUgras{võtmeid} ja laulma kauneid \VUgras{meloodiaid}. Noodid pannakse
\VUgras{noodipuldile} \textsuperscript{(65)}. Üks õpilastest \VUgras{mängib klaverit}
\textsuperscript{(64)} ja \VUgras{muusikaõpetaja} \textsuperscript{(66)} \VUgras{lööb
takti} \textsuperscript{(68)}.

%%K t01-17-1 p
17. --- Me hakkame nägema nurka \VUgras{käsitöö} \VUgras{töökojast}
\textsuperscript{(69)}, kus poisid võivad õppida puitu hööveldama ja rauda viilima.
Seda ei ole igas gümnaasiumis.

%%K t01-tit-7 sub
\VUsaut{5.0mm}
\VUpk{10.2pt}{40}{Loodusteaduste klass.}
\VUsaut{3.6mm}

%%K t01-18-1 p
18. --- Matemaatikaõpetaja õpetab õpilastele \VUgras{geomeetriat, aritmeetikat,
arvutamist} ja \VUgras{meetermõõdustikku}. Me näeme tabelil peamisi geomeetrilisi
kujundeid: \VUgras{ring} \textsuperscript{(a)}, \VUgras{ruut} \textsuperscript{(b)},
\VUgras{kolmnurk} \textsuperscript{(c)}, \VUgras{joon} \textsuperscript{(d)}.

%%K t01-18-2 p
Me näeme ka üht \VUgras{tehet}; see ei ole \VUgras{liitmine} ega \VUgras{lahutamine}
ega \VUgras{korrutamine}: see on \VUgras{jagamine} \textsuperscript{(e)}.

%%K t01-19-1 p
19. --- Meetermõõdustiku tabelil eristame: \VUgras{meetrit} \textsuperscript{(a)},
\VUgras{kaalu} \textsuperscript{(b)} ja selle \VUgras{vihte}, \VUgras{liitrit}
\textsuperscript{(e)}, \VUgras{dekaliitrit} \textsuperscript{(f)}, \VUgras{detsiliitrit}
ja \VUgras{kuupmeetrit} \textsuperscript{(d)}.

%%K t01-19-2 p
Õpilased õpivad ka \VUgras{loodusteadusi} \textsuperscript{(11)} --- zooloogiat
\textsuperscript{(a)}, botaanikat \textsuperscript{(b)}, geoloogiat
\textsuperscript{(c)} --- ja \VUgras{ajalugu} \textsuperscript{(12)}.

%%K t01-19-3 p
Kapis on \VUgras{füüsika} \textsuperscript{(15)} ja \VUgras{keemia}
\textsuperscript{(16)} riistad.

%%K t01-19-4 p
Kapi peal on: \VUgras{kera} \textsuperscript{(14)}, \VUgras{koonus} ja
\VUgras{maakera} \textsuperscript{(13)}.

%%K t01-19-5 p
Tahvlite kohal ripub \VUgras{kaart}, mis kujutab maailma viit jagu:
\VUgras{Euroopat} \textsuperscript{(g)}, \VUgras{Aasiat} \textsuperscript{(e)},
\VUgras{Aafrikat} \textsuperscript{(f)}, \VUgras{Ameerikat} \textsuperscript{(ab)} ja
\VUgras{Okeaaniat} \textsuperscript{(h)}, ning kaht suurt ookeani, \VUgras{Vaikset}
\textsuperscript{(c)} ja \VUgras{Atlandi} \textsuperscript{(d)}.
\VUsaut{6.0mm}
\VUfilet{22mm}
